\section{The original coefficient flags and the first heat response of the socle}
\label{sec:heat-flag-response}

This calculation starts with the actual heat family, coordinates and
full analytic unit of HC1--4 and HH1--6. Those proofs retain the
Rodgers--Tao convention $H_t(z)=\int_0^\infty e^{tu^2}\Phi(u)\cos(zu)\,du$,
$\partial_tH_t=-\partial_z^2H_t$, and
$H_0(z)=\xi_{\mathrm R}(1/2+iz/2)/8$ from
Brad Rodgers and Terence Tao,
\href{https://arxiv.org/abs/1801.05914v5}{\emph{The de Bruijn--Newman constant is non-negative}},
author equations \texttt{hoz}, \texttt{htdef} and \texttt{sas}.
The original local Weil contribution and its full multiplicity are
identified in HH23--30 using Alain Connes and Caterina Consani,
\href{https://arxiv.org/abs/2207.10419}{\emph{Hochschild homology, trace map and zeta-cycles}},
Proposition \texttt{laplacian} and Remark \texttt{remtauto}, and
\href{https://arxiv.org/abs/2006.13771v1}{\emph{Weil positivity and trace formula, the archimedean place}}.
All finite algebra, limits and coefficients used below are proved here.

Fix an actual real zero $x_0$ of order $m\ge1$. Write $\xi=z-x_0$,
$t=u^2$, and retain the entire original polynomial
$W(t,\xi)=\prod_{\nu=1}^m(\xi-\theta_\nu(u))$ and unit
$H_t(x_0+\xi)=U(t,\xi)W(t,\xi)$. Put
\[
 c=a_{m+1}/a_m,\quad d=a_{m+2}/a_m,\quad b=2d-c^2,
 \qquad
 \theta_\nu(u)=r_\nu u+2cu^2+b r_\nu u^3+O(u^4).
 \tag{HF1}
\]
The $r_\nu$ are all the distinct real roots of
\[
 P_m(y)=\sum_{a=0}^{\lfloor m/2\rfloor}
 \frac{(-1)^a m!}{a!(m-2a)!}y^{m-2a}.
 \tag{HF2}
\]
No multiple zero is asserted to exist. The construction applies to
each actual zero with its actual order. In particular it includes $m=1$.

\paragraph{HF1. The full original trace matrix.}
Let $\mathcal O=\mathbb C\{t\}$,
$A=\mathcal O[\xi]/(W)$ and retain its ordered basis
$1,\xi,\ldots,\xi^{m-1}$. For real $t$, conjugation fixes $t,\xi$;
write $f^\#(\xi)=\overline{f(\bar\xi)}$. The reflected regular trace is
\[
 T_t(f,g)=\operatorname{Tr}_{A_t/\mathbb C}(M_{f^\#g}),
 \qquad K_{ij}(t)=s_{i+j}(t),\qquad
 s_n(t)=\sum_{\nu=1}^m\theta_\nu(u)^n,
 \quad 0\le i,j<m.
 \tag{HF3}
\]
The sum is unchanged by $u\mapsto-u$, which permutes the branches;
it is thus analytic in $t$. At $t\ne0$ the algebra is reduced and
the formula is evaluation of its multiplication trace. At zero it
continues to that same trace by the companion matrix of $W$.
Consequently HF3 is an exact identity at every fibre, including
the nonreduced one, and $K(0)=\operatorname{diag}(m,0,\ldots,0)$.

For $0\le j<m$ define the original coefficient submodule
\[
 F^jA=\operatorname{span}_{\mathcal O}
          \{\xi^j,\xi^{j+1},\ldots,\xi^{m-1}\}.
 \tag{HF4}
\]
These are specified free submodules in the original basis. They
are not silently identified with ideals of all deformed fibres.
Their specializations are the ideals $(\xi^j)$ of
$A_0=\mathbb C[\xi]/(\xi^m)$.

\paragraph{HF2. All successive trace responses are evaluated.}
Put $p_n=\sum_\nu r_\nu^n$, including $p_0=m$. For $f,g\in F^jA$
with coefficient functions $f_k(t),g_k(t)$, the exact limit is
\[
 \boxed{\lim_{t\to0}t^{-j}T_t(f(t),g(t))
       =p_{2j}\,\overline{f_j(0)}g_j(0).}
 \tag{HF5}
\]
Here the limit is along real $t$; the matrix quotient has an analytic
continuation through zero. To prove it, HF1 gives
$s_{2a}(t)=p_{2a}t^a+O(t^{a+1})$, while evenness in $u$ gives
$s_{2a+1}(t)=O(t^{a+1})$. If $i,j\ge j_0$ and
$i+j>2j_0$, then $s_{i+j}(t)=O(t^{j_0+1})$.
The only entry of $t^{-j_0}K|_{F^{j_0}}$ surviving at zero is
therefore $(j_0,j_0)$, with value $p_{2j_0}$. Multiplying by the
analytic coefficient columns proves HF5.

For $j=0$, $p_0=m>0$. For $m\ge2$ and $j>0$, at least one
$r_\nu$ is nonzero, so $p_{2j}>0$. Thus HF5 has radical
$F^{j+1}A_0$ inside $F^jA_0$ and is positive definite on the
one-dimensional quotient $F^jA_0/F^{j+1}A_0$.
Every endpoint nilpotent layer has an explicitly determined
first heat response. This is a statement about the time derivatives
of the retained family, not a new positive value assigned to $T_0$.
In particular,
\[
 \left.\frac d{dt}T_t(f,g)\right|_{t=0}
    =2m(m-1)\,\overline{f_1(0)}g_1(0),
       \qquad f,g\in F^1A.
 \tag{HF6}
\]
The identity $p_2=2m(m-1)$ follows from the coefficient
$-m(m-1)$ of $y^{m-2}$ in HF2 and Newton's second identity.

\paragraph{HF3. The coefficients are finite exact data.}
Writing $P_m(y)=y^m+\alpha_1y^{m-1}+\cdots+\alpha_m$, all
$\alpha_{2a}=(-1)^a m!/(a!(m-2a)!)$ and all odd $\alpha_k$ vanish.
The required $p_n$, including $p_{2m-2}$, are evaluated by
\[
 \begin{cases}
 p_n+\alpha_1p_{n-1}+\cdots+\alpha_{n-1}p_1+n\alpha_n=0,
          &1\le n\le m,\\
 p_n+\alpha_1p_{n-1}+\cdots+\alpha_m p_{n-m}=0,&n>m.
 \end{cases}
 \tag{HF7}
\]
For completeness, expand
$P_m'(y)/P_m(y)=\sum_\nu(y-r_\nu)^{-1}
=\sum_{n\ge0}p_n y^{-n-1}$ at infinity, multiply by $P_m(y)$,
and equate coefficients. This proves HF7 without numerical root
approximation. The Hermite convention in HF2 agrees with
\href{https://dlmf.nist.gov/18.5.E13}{DLMF18.5.13}, Chapter18 by
T.~H. Koornwinder, R.~Wong, R.~Koekoek and R.~F. Swarttouw;
the needed polynomial and Newton calculation are explicit above.

The next coefficient still retains the actual arithmetic unit data.
For each $j\ge1$ one has
\[
 \boxed{s_{2j}(t)=t^j\left[
    p_{2j}+t\left(2jb\,p_{2j}
                +4j(2j-1)c^2p_{2j-2}\right)+O(t^2)\right].}
 \tag{HF8}
\]
Indeed the terms of degree $u^{2j+2}$ in HF1 raised to the
$2j$-th power are $2jb r_\nu^{2j}$ and
$\binom{2j}{2}(2c)^2r_\nu^{2j-2}$. The possible degree
$2j+1$ coefficient sums to zero by symmetry of the $r_\nu$.
Every unlisted term has higher degree, and the full sum is even
in $u$, proving the $O(t^2)$ remainder inside brackets.
No value of $c$ or $d$ is replaced by that of a model polynomial.

\paragraph{HF4. The original trace-invisible socle has a measured response.}
The arithmetic coordinate is $v=s-\rho=i\xi/2$. For a fixed
complex amplitude $\beta$, the original socle class
$\beta v^{m-1}$ has the specified coefficient lift
\[
 f_\beta(t,\xi)=\beta(i\xi/2)^{m-1}\in A.
 \tag{HF9}
\]
For $m\ge2$ its value under $T_0$ is zero. Its complete deformed
trace and first two nonzero orders are
\[
 \begin{split}
 T_t(f_\beta,f_\beta)
   &=\frac{|\beta|^2}{4^{m-1}}s_{2m-2}(t)\\
   &=\frac{|\beta|^2t^{m-1}}{4^{m-1}}
       \left[p_{2m-2}+t\left(
           2(m-1)b p_{2m-2}
          +4(m-1)(2m-3)c^2p_{2m-4}\right)+O(t^2)\right].
 \end{split}
 \tag{HF10}
\]
The first equality follows directly from HF3 and complex conjugation
of the factor $i^{m-1}$; the second is HF8. For nonzero $\beta$
the leading coefficient is strictly positive. For small negative
$t$ its sign is $(-1)^{m-1}$; for small positive $t$ it is positive.
These signs concern the explicitly deformed trace and the specified
lift. They do not replace the zero arithmetic endpoint contribution.
For $m=2$, the exact first two orders are
\[
 T_t(f_\beta,f_\beta)
   =|\beta|^2\bigl(t+4d\,t^2+O(t^3)\bigr),
 \tag{HF11}
\]
because $p_2=4$, $p_0=2$ and $b+c^2=2d$.
For $m=3$, HF7 gives $p_2=12$, $p_4=72$, and hence
$T_t(f_\beta,f_\beta)=|\beta|^2t^2(9/2+36dt+O(t^2))$.
For $m=4$, it gives $p_4=240$, $p_6=2592$, and hence
$T_t(f_\beta,f_\beta)=|\beta|^2t^3(81/2+(486d-18c^2)t+O(t^2))$.
These formulas evaluate the full constants at each displayed order;
they do not assert the existence of those multiplicities.

\paragraph{HF5. A finite algebra retaining the whole separated boundary.}
The weighted trace response has an exact algebra map. Over
$\mathcal O_u=\mathbb C\{u\}$ put
\[
 \widetilde W(u,y)=u^{-m}W(u^2,uy),\qquad
 \mathscr R=\mathcal O_u[y]/(\widetilde W),\qquad
 A\otimes_{\mathcal O}\mathcal O_u\longrightarrow\mathscr R,
       \quad\xi\longmapsto uy.
 \tag{HF12}
\]
Every coefficient of $\widetilde W$ is analytic: its roots are
$z_\nu(u)=\theta_\nu(u)/u$, which extend analytically to
$z_\nu(0)=r_\nu$. All their pairwise differences are units at zero.
Thus evaluation gives an explicit isomorphism
$\mathscr R\simeq\prod_{\nu=1}^m\mathcal O_u$, with inverse
the Lagrange polynomials
$\prod_{\mu\ne\nu}(y-z_\mu)/(z_\nu-z_\mu)$.
This is not a replacement of $A$; HF12 keeps its injection and
all coefficients of its defining polynomial. In the original
ordered power bases the map is exactly
$\operatorname{diag}(1,u,\ldots,u^{m-1})$. Hence
\[
 \mathscr R/(A\otimes\mathcal O_u)
     \simeq\bigoplus_{j=0}^{m-1}\mathcal O_u/(u^j),\qquad
 \operatorname{length}=\frac{m(m-1)}2.
 \tag{HF13}
\]
For real nonzero $u$, if $\widetilde K_{ij}(u)=\sum_\nu z_\nu(u)^{i+j}$
and $D_u=\operatorname{diag}(1,u,\ldots,u^{m-1})$, the exact
metric identity is
\[
 K(u^2)=D_u^*\widetilde K(u)D_u,\qquad
 \widetilde K(0)=(p_{i+j})_{0\le i,j<m},\qquad
 \det\widetilde K(0)=2^{m(m-1)/2}\prod_{j=1}^m j^j>0.
 \tag{HF14}
\]
The matrix at zero is the Gram matrix of evaluation at all
distinct real $r_\nu$, hence positive definite. Its determinant
is the squared Vandermonde, evaluated in HC14. The original
Gram matrix remains $K(0)$; the exact factors $D_u$ record its
loss of rank. This also proves that all flag limits HF5 belong
to one full family, rather than to unrelated rescalings.

The intermediate algebra of HC5 supplies the commutative diagram
\[
 \begin{tikzcd}
 A\otimes\mathcal O_u\arrow[r,hook]\arrow[dr,hook]&
 B\otimes\mathcal O_u\arrow[d,hook]\\
 &\mathscr R.
 \end{tikzcd}
 \tag{HF15}
\]
On each paired component the right arrow is
$\mathcal O_u[v]/(v^2-u^2)\to\mathcal O_u^2$,
$v\mapsto(u,-u)$. Its matrix in bases $(1,v)$ and the two
idempotents is $\left(\begin{smallmatrix}1&u\\1&-u\end{smallmatrix}\right)$,
so its cokernel is $\mathcal O_u/(u)$. The central component,
if present, maps identically. With $r=\lfloor m/2\rfloor$ this proves
\[
 \operatorname{length}\bigl(\mathscr R/(B\otimes\mathcal O_u)\bigr)=r,
 \qquad \frac{m(m-1)}2
     =2\left\lfloor\frac{(m-1)^2}{4}\right\rfloor+r.
 \tag{HF16}
\]
Flatness of $\mathcal O_u$ over $\mathcal O$ follows from its free
basis $(1,u)$, so the lower index is exactly twice HC8's index.
Finally, applying $G_L$ to every injective algebra map in HF15
gives the corresponding injective semiring maps, with
$(a,1_L)\mapsto(\phi(a),1_L)$ and $(0,\lambda)\mapsto(0,\lambda)$.
Injectivity ensures that a nonzero amplitude never becomes a zero
label. Addition and multiplication commute with these maps
componentwise. The entire original support lattice and both
amplitude/support projections therefore survive this boundary
construction.

\begin{figure}[p]
\centering
\begin{tikzpicture}[>=Stealth,
 box/.style={draw,rounded corners,align=center,text width=11cm,inner sep=8pt}]
\node[box] (a) at (0,0)
 {Original coefficient flag $F^jA$\\
 $\xi^j,\xi^{j+1},\ldots,\xi^{m-1}$; full $W$ and $U$ retained};
\node[box] (b) at (0,-2.3)
 {$t^{-j}T_t|_{F^j}$ extends to the endpoint\\
 $p_{2j}\,\overline{f_j(0)}g_j(0)$,
 \quad radical $F^{j+1}A_0$};
\draw[->] (a)--node[right,align=left]
 {exact original trace\\HF3--5}(b);
\node[box] (c) at (0,-4.6)
 {$F^jA_0/F^{j+1}A_0$ has weight $p_{2j}>0$\\
 Each successive nilpotent layer is detected at its own time order.};
\draw[->] (b)--node[right]{induced one-dimensional pairing}(c);
\end{tikzpicture}
\par\vspace{6mm}
\renewcommand{\arraystretch}{1.4}
\begin{tabular}{ccl}
\toprule
Multiplicity&First time order&$T_t(\beta v^{m-1},\beta v^{m-1})/|\beta|^2$\\
\midrule
$2$&$t$&$t+4dt^2+O(t^3)$\\
$3$&$t^2$&$\frac92t^2+36dt^3+O(t^4)$\\
$4$&$t^3$&$\frac{81}2t^3+(486d-18c^2)t^4+O(t^5)$\\
\bottomrule
\end{tabular}
\caption{The original endpoint nilpotent flag and its evaluated heat
responses, HF1--11. The coordinate remains $v=i\xi/2$ and the
arithmetic coefficients remain $c=a_{m+1}/a_m$, $d=a_{m+2}/a_m$.
The table describes the exact formulas for those orders, without
asserting that a multiple zeta zero exists. The source heat convention
is Rodgers--Tao v5; the original reflected arithmetic trace is the
Connes--Consani receiver cited at the start of this section.}
\end{figure}
